% ---------------------------------------------------------------------
%  2. Estado del arte
% ---------------------------------------------------------------------
\section{Estado del arte}
\label{sec:estado-arte}

\subsection{El problema: teoría de la información y codificación sin pérdida}
\label{sec:ea-problema}
El punto de partida de toda la disciplina es el trabajo de Shannon~\cite{shannon1948}, que
modela una fuente como un proceso aleatorio y define su entropía como el número mínimo
promedio de bits por símbolo con que puede representarse sin pérdida. El teorema de
codificación de fuente convierte la compresión en un problema bien planteado: existe un límite
inferior y la calidad de un compresor se mide por su distancia a él. Sin embargo, la entropía
depende del \emph{modelo} de la fuente: un modelo de orden cero, que trata cada byte como
independiente, ignora la redundancia contenida en secuencias repetidas, mientras que modelos
con contexto capturan dependencias entre símbolos consecutivos~\cite{bell1989}.

Bell, Witten y Cleary~\cite{bell1989} propusieron una distinción que sigue organizando el campo:
todo compresor combina un \emph{modelo} que asigna probabilidades o identifica repeticiones y
un \emph{codificador} que traduce esa información en bits. Bajo esa óptica, los métodos se
agrupan en \emph{estadísticos}, que codifican símbolos según su probabilidad, y \emph{de
diccionario}, que sustituyen secuencias por referencias. Los textos de referencia de
Sayood~\cite{sayood2017} y de Salomon y Motta~\cite{salomon2010} adoptan esta misma taxonomía, y
la revisión de Jayasankar et al.~\cite{jayasankar2021} la amplía al clasificar las técnicas por
calidad de reconstrucción, esquema de codificación, tipo de dato y aplicación. Una conclusión
común a estas revisiones es que no existe un compresor universalmente mejor: el desempeño
depende de la coincidencia entre la redundancia que explota el método y la que contienen los
datos. Esta observación motiva comparar Huffman y LZW sobre fuentes de naturaleza
deliberadamente distinta.

\subsection{Algoritmos que resuelven el problema}
\label{sec:ea-algoritmos}
\paragraph{Familia de Huffman.} Huffman~\cite{huffman1952} demostró que un algoritmo voraz que
fusiona repetidamente los dos símbolos menos probables produce un código prefijo de longitud
media mínima para una distribución conocida. Gallager~\cite{gallager1978} caracterizó estos
códigos mediante la propiedad de hermanos (\emph{sibling property}), acotó su redundancia en
función de la probabilidad del símbolo más frecuente y planteó una versión adaptativa. Esa
línea se completó con los algoritmos dinámicos de Knuth~\cite{knuth1985} y de
Vitter~\cite{vitter1987}, que actualizan el árbol tras cada símbolo y evitan la segunda pasada
y la transmisión de la tabla, a cambio de un mayor costo por símbolo. Desde el punto de vista
de la implementación, Moffat y Turpin~\cite{moffat1997} mostraron que basta transmitir las
longitudes de los códigos si se usan \emph{códigos canónicos}, que además permiten decodificar
con tablas pequeñas; la revisión de Moffat~\cite{moffat2019} sintetiza siete décadas de estas
mejoras. La limitación estructural de toda la familia es que asigna un número entero de bits a
cada símbolo, por lo que pierde hasta un bit por símbolo frente a la entropía cuando la
distribución está muy sesgada.

\paragraph{Codificación aritmética: alternativa fuera de alcance.} La codificación
aritmética~\cite{witten1987} supera esa limitación al representar el mensaje completo como un
intervalo y acercarse a la entropía del modelo con una fracción de bit de sobrecosto por
mensaje. Constituye el estándar en compresores estadísticos avanzados, pero su comparación
excede el alcance de este trabajo, centrado en los dos algoritmos clásicos de cada familia.

\paragraph{Familia Lempel--Ziv.} Ziv y Lempel propusieron dos esquemas universales que no
requieren conocer la distribución de la fuente: LZ77~\cite{ziv1977}, que codifica cada frase
como una referencia a una coincidencia previa dentro de una ventana deslizante, y
LZ78~\cite{ziv1978}, que construye incrementalmente un diccionario explícito de frases y cuya
tasa converge asintóticamente a la entropía para fuentes estacionarias y ergódicas. Storer y
Szymanski~\cite{storer1982} formalizaron la compresión por sustitución textual y propusieron
evitar referencias que no ahorran espacio, idea conocida como LZSS. Welch~\cite{welch1984}
simplificó LZ78 inicializando el diccionario con todos los símbolos del alfabeto, de modo que
la salida consiste solo en índices; el resultado, LZW, es rápido y sencillo de implementar
en hardware y software, pero su diccionario crece sin límite si no se gestiona.

\paragraph{DEFLATE: la combinación de ambas ideas.} El formato DEFLATE~\cite{rfc1951} combina
una etapa LZ77 con una codificación de Huffman canónica de los literales, longitudes y
distancias resultantes. Su éxito práctico ilustra que ambas familias son complementarias: la
etapa de diccionario elimina la repetición de frases y la etapa estadística, el sesgo en la
distribución de los símbolos resultantes. DEFLATE queda fuera del alcance de este trabajo, que
compara los dos mecanismos por separado.

\paragraph{Gestión del diccionario en LZW.} Cuando el diccionario de LZW se llena, el
compresor debe decidir qué hacer. Cheng et al.~\cite{cheng2021} comparan las políticas
habituales: el reinicio completo (\emph{reset}) es más rápido pero empeora la tasa de
compresión, mientras que el reemplazo LRU obtiene mejor tasa a costa de más tiempo y memoria;
proponen IDA-LZW, que gestiona un área de entradas poco significativas. La implementación de
este trabajo adopta el reinicio, la política más simple y predecible, y su efecto se estudia
variando el tamaño máximo del diccionario.

\subsection{Comparaciones previas y metodologías de evaluación}
\label{sec:ea-comparaciones}
La \cref{tab:trabajos-relacionados} resume los estudios que comparan directamente Huffman y
LZW o variantes de ambos. Jambek y Khairi~\cite{jambek2014} los evalúan para nodos de redes de
sensores inalámbricos y reportan que Huffman redujo los datos en alrededor de 43\,\% en
promedio y fue aproximadamente cuatro veces más rápido que LZW. Fauzan et
al.~\cite{fauzan2022}, en cambio, trabajan con archivos de texto (\texttt{.txt}, \texttt{.csv},
entre otros) y encuentran que LZW obtuvo mayor ahorro de espacio en \texttt{.txt} y
\texttt{.csv} (63.85\,\% y 77.56\,\% en promedio), con tiempos de compresión que crecen con
el tamaño del archivo. Btoush et al.~\cite{btoush2008} comparan LZW, Huffman y codificación de
longitud fija sobre textos de longitud variable, midiendo tasa, tiempo y entropía, y observan
que LZW es el peor en archivos pequeños. Bjärås~\cite{bjaras2019} evalúa cinco algoritmos
ligeros para series temporales en dispositivos IoT, incluidas variantes basadas en Huffman y en
LZW, y concluye que las basadas en Huffman ofrecen el mejor equilibrio entre tasa, memoria y
uso de CPU. Cheng et al.~\cite{cheng2021} no comparan contra Huffman, pero muestran que la
política de diccionario lleno cambia por sí sola la tasa y el tiempo de LZW.

Estos resultados, aparentemente contradictorios, son coherentes con la observación de la
\cref{sec:ea-problema}: las lecturas de sensores y las series temporales tienen poca
repetición de frases y una distribución de valores sesgada, lo que favorece a Huffman,
mientras que el texto contiene palabras y frases recurrentes, que LZW aprovecha una vez que su
diccionario se ha poblado; en archivos pequeños el diccionario no alcanza a aprender, lo que
explica la desventaja reportada por Btoush et al. Ninguno de estos trabajos, sin embargo,
aísla experimentalmente cada tipo de redundancia ni verifica cómo escalan los tiempos.

Respecto a la metodología, la literatura de algoritmia experimental ofrece pautas que los
estudios anteriores no siguen. Johnson~\cite{johnson2002} recomienda instancias de prueba
diversas y de tamaño creciente, informes reproducibles y comparaciones con la teoría.
McGeoch~\cite{mcgeoch2012} describe cómo estimar la tasa de crecimiento de un algoritmo a
partir de mediciones, por ejemplo mediante la pendiente en escala log--log. Georges et
al.~\cite{georges2007} demuestran que, en Java, la compilación JIT, la recolección de basura y
la variabilidad entre ejecuciones pueden invertir las conclusiones si se reporta una única
medición, y proponen repeticiones con intervalos de confianza. Por último, el uso de corpus
públicos como Canterbury~\cite{arnold1997} y Silesia~\cite{deorowicz2003} hace comparables los
resultados entre estudios.

\begin{table}[H]
  \centering
  \caption{Estudios comparativos relacionados.}
  \label{tab:trabajos-relacionados}
  \footnotesize
  \setlength{\tabcolsep}{4pt}
  \renewcommand{\arraystretch}{1.15}
  \begin{tabularx}{\linewidth}{@{}>{\raggedright}p{1.6cm} >{\raggedright}p{2.4cm} >{\raggedright}p{2.6cm}
                  >{\raggedright}p{2.2cm} >{\raggedright\arraybackslash}X@{}}
    \toprule
    Trabajo & Algoritmos & Datos & Métricas & Hallazgo principal \\
    \midrule
    \cite{jambek2014} & Huffman, LZW & Datos de nodos de sensores inalámbricos &
      Reducción de datos, tiempo &
      Huffman redujo los datos alrededor de 43\,\% en promedio y fue unas cuatro veces más
      rápido que LZW. \\
    \cite{fauzan2022} & Huffman, LZW & Archivos de texto (\texttt{.txt}, \texttt{.csv}, otros) &
      Ahorro de espacio, tiempo &
      LZW logró mayor ahorro en \texttt{.txt} y \texttt{.csv} (63.85\,\% y 77.56\,\% en
      promedio); el tiempo creció con el tamaño del archivo. \\
    \cite{btoush2008} & LZW, Huffman, longitud fija & Textos de longitud variable &
      Tasa, tiempo, entropía & LZW fue el peor en archivos pequeños. \\
    \cite{bjaras2019} & Cinco algoritmos ligeros (variantes Huffman y LZW, entre otros) &
      Series temporales IoT & Tasa, memoria, CPU &
      Las variantes basadas en Huffman resultaron las más equilibradas. \\
    \cite{cheng2021} & Políticas de diccionario de LZW; IDA-LZW & Conjuntos de prueba propios del estudio &
      Tasa, tiempo, memoria &
      El reinicio es más rápido pero comprime peor; LRU comprime mejor con más tiempo y
      memoria. \\
    \midrule
    Este trabajo & Huffman canónico, LZW 9--16 bits &
      Sintéticos controlados, texto real, Canterbury y Silesia &
      Tiempo (mediana, IC 95\,\%), pendiente log--log, bpb &
      Ver \cref{sec:resultados,sec:discusion}. \\
    \bottomrule
  \end{tabularx}
\end{table}

\subsection{Síntesis y brecha}
\label{sec:ea-brecha}
La revisión permite identificar tres carencias en las comparaciones previas entre Huffman y
LZW. Primero, los datos son pequeños o de un único dominio (sensores, series temporales o
texto), lo que impide separar el efecto del tipo de redundancia del efecto del algoritmo.
Segundo, ningún estudio valida la complejidad teórica: los tiempos se reportan para unos pocos
tamaños, sin estimar la tasa de crecimiento. Tercero, falta rigor estadístico: no se describen
calentamiento, repeticiones ni intervalos de confianza, pese a que en entornos con JIT esto
puede alterar las conclusiones~\cite{georges2007}.

Este trabajo aborda esas carencias con (i)~fuentes sintéticas que aíslan cada tipo de
redundancia, junto con texto real y los corpus estándar~\cite{arnold1997,deorowicz2003};
(ii)~tamaños de entrada que crecen geométricamente, de $2^{14}$ a $2^{24}$ bytes, y un
ajuste log--log contrastado con la complejidad derivada~\cite{mcgeoch2012}; y (iii)~un
protocolo con calentamiento, repeticiones, medianas e intervalos de confianza de Student,
siguiendo las recomendaciones de la algoritmia experimental~\cite{johnson2002,georges2007}.
